% 図：tmux で画面を分割した様子（chapters/tools/network.tex）
\begin{tikzpicture}[pane/.style={draw, fill=black!3, anchor=north west, inner sep=2pt, align=left, font=\scriptsize\ttfamily}]
  \def\W{9.0} \def\H{4.2}
  \fill[ubuntutitlebar] (0,0) rectangle (\W,0.45);
  \node[font=\footnotesize\bfseries, text=white] at (\W/2,0.225) {user@robot: \textasciitilde};
  \draw[ubuntutitlebar, thick] (0,0.45) rectangle (\W,-\H);
  \draw[thick] (\W/2,0) -- (\W/2,-\H+0.35);
  \draw[thick] (\W/2,-1.9) -- (\W,-1.9);
  \node[pane, draw=none, fill=none] at (0.05,-0.05) {\$ ros2 launch my\_robot\\\ \ bringup.launch.py\\{[INFO]} [robot\_driver]: started\\{[INFO]} [lidar\_node]: started\\...};
  \node[pane, draw=none, fill=none] at (\W/2+0.05,-0.05) {\$ ros2 topic list\\/cmd\_vel\\/scan\\/odom};
  \node[pane, draw=none, fill=none] at (\W/2+0.05,-1.95) {\$ htop\\CPU [|||||      35\%]\\Mem [||||     1.2G]};
  \fill[green!50!black] (0,-\H+0.35) rectangle (\W,-\H);
  \node[font=\scriptsize\ttfamily, text=white, anchor=west] at (0.05,-\H+0.175) {[robot] 0:bash*};
  \node[figlabel, anchor=north, text=black!60] at (\W/2,-\H-0.1) {1 つのターミナルの中で、3 つのシェルを同時に使っている};
\end{tikzpicture}
